% mode: ini
% Aimed at hyphenate (sections 1072-1095) as it meets ordinary text with real fonts: Knuth's patterns from hyphen.tex, words with fi, ff, ffi, fl ligatures and kerns that must be rebuilt in the reconstitution (section 1085's init_list, 1088's
% ligature replacement, section 1092's pre_break with the hyphen, 1093's post_break, 1094's catching up, the replace count of 1095); \lefthyphenmin, \righthyphenmin; a font whose hyphen char is absent (new_character's message, hyf_node=null), \hyphenchar=-1 (no try); \showbox at full depth shows the result node by node.
\input prelude
\lccode`\'=`\'
\input hyphen
\lefthyphenmin=2 \righthyphenmin=3 \defaulthyphenchar=`\- \hyphenpenalty=-1000 \exhyphenpenalty=50 \hsize=1pt \tolerance=10000 \parfillskip=0pt \parindent=0pt \pretolerance=-1 \linepenalty=10
\tenrm \hyphenchar\tenrm=`\-
\tracingcommands=0 \tracingonline=1
% Twice each: in a narrow measure, so that each hyphen breaks a line, and in a wide one, where the discretionaries stay whole and show pre_break, post_break and "replacing".
\def\try#1{{\hsize=1pt \setbox0=\vbox{\hskip0pt #1\par}\showbox0 }{\hsize=1000pt \hyphenpenalty=50 \linepenalty=10 \tolerance=-1 \setbox0=\vbox{\hskip0pt #1\par}\showbox0 }}
\try{hyphenation}
\try{difficult}
\try{affluent}
\try{office}
\try{coffee}
\try{flying}
\try{effervescent}
\try{sufficient}
\try{waffle}
\try{awkward}
\try{typewriter}
\try{attenuation}
\try{ligature}
\try{AVAST}
\try{maxwell}
\try{associate}
\try{associates}
\try{table}
\try{oxygen}
\try{encyclopaedia}
% Letters and the mins.
\lefthyphenmin=1 \righthyphenmin=1
\try{hyphenation}
\try{office}
\lefthyphenmin=5 \righthyphenmin=5
\try{hyphenation}
\try{office}
\lefthyphenmin=32 \righthyphenmin=32
\try{hyphenation}
\lefthyphenmin=63 \righthyphenmin=64
\try{hyphenation}
\lefthyphenmin=0 \righthyphenmin=0
\try{hyphenation}
\lefthyphenmin=-1 \righthyphenmin=-1
\try{hyphenation}
\lefthyphenmin=2 \righthyphenmin=3
% Word delimiters: what starts and ends a word.
\try{hyphenation,}
\try{hyphenation.}
\try{(hyphenation)}
\try{hyphenation's}
\try{hyphenation\/}
\try{hyphenation\kern 1pt}
\try{hyphenation\hskip 1pt}
\try{hyphenation\penalty 0}
\try{hyphenation\discretionary{}{}{}}
\try{hyphenation\vrule}
\try{hyphenation\mark{m}}
\try{hyphenation\special{s}}
\try{hyphenation\insert100{}}
\try{hyphenation\vadjust{}}
\try{hyphenation$ $}
\try{hyphenation\hbox{}}
\try{hyphen\/ation}
\try{hyphen\kern 0pt ation}
\try{hyphen\-ation}
\try{hy\-phen\-ation}
\try{hyphen\accent 18 ation}
\try{\hbox{}hyphenation}
\try{\mark{m}hyphenation}
\try{\kern 1pt hyphenation}
% A word over 63 letters is not touched; 63 exactly is.
\try{aaaaaaaaaaaaaaaaaaaaaaaaaaaaaaaaaaaaaaaaaaaaaaaaaaaaaaaaaaaaaaaaaaaaaation}
\try{hyphenationhyphenationhyphenationhyphenationhyphenationhyphena}
\try{hyphenationhyphenationhyphenationhyphenationhyphenationhyphenat}
\try{hyphenationhyphenationhyphenationhyphenationhyphenationhyphenati}
% The hyphen char: absent, 256 or -1.
\hyphenchar\tenrm=`\A
\try{hyphenation}
\hyphenchar\tenrm=`\-
\font\xr=cmr10 at 12pt \xr
\hyphenchar\xr=127
\try{hyphenation}
\hyphenchar\xr=128
\try{hyphenation}
\tracinglostchars=1
\hyphenchar\xr=200
\try{hyphenation}
\tracinglostchars=0
\hyphenchar\xr=255
\try{hyphenation}
\hyphenchar\xr=256
\try{hyphenation}
\hyphenchar\xr=-1
\try{hyphenation}
\hyphenchar\xr=0
\try{hyphenation}
\tenrm
\try{Hyphenation}
\uchyph=0
\try{Hyphenation}
\uchyph=1
% A font with no ligatures or kerns, and a typewriter.
\font\tt=cmtt10 \tt \hyphenchar\tt=`\-
\try{hyphenation}
\try{difficult}
\hyphenchar\tt=-1
\try{hyphenation}
\tenrm
\try{{\tt hyphen}ation}
\try{hyphen{\tt ation}}
\end
